\chapter{Efim Zelmanov --- Fields Medal 1994}\label{ch1994efimzelmanov:efim-zelmanov-fields-medal-1994-a-comprehensive-report-on-his-work-and-impact}

\begin{center}\rule{0.5\linewidth}{0.5pt}\end{center}

\textbf{Prepared:} July 2026 \textbf{Subject:} Efim Isaakovich Zelmanov
(b. September 18, 1955, Khabarovsk, USSR) \textbf{Position:} Professor,
University of California, San Diego (since 2002) \textbf{Doctoral
Advisor:} Anatoly Shirshov (Novosibirsk State University) \textbf{Fields
Medal Citation:} \emph{``For his solution of the restricted Burnside
problem and his contributions to the theory of Jordan algebras and Lie
algebras.''}

\begin{center}\rule{0.5\linewidth}{0.5pt}\end{center}

\section{1. Introduction}\label{ch1994efimzelmanov:introduction}

Efim Zelmanov was awarded the Fields Medal in 1994 for his landmark
solution of the \textbf{restricted Burnside problem}, a fundamental
problem in group theory that had remained open for over 90 years. His
work also includes significant contributions to \textbf{Jordan algebras}
and \textbf{Lie algebras}, including the classification of prime Jordan
algebras and the solution of the Kurosh problem.

Zelmanov's work represents a triumph of the structural approach to
algebra. By developing deep structural theories of non-associative
algebras and applying them to group theory, he solved one of the most
famous problems in mathematics.

This report provides a comprehensive account of Zelmanov's major
contributions.

\begin{center}\rule{0.5\linewidth}{0.5pt}\end{center}

\section{2. Biographical Background}\label{ch1994efimzelmanov:biographical-background}

\subsection{2.1 Early Life}\label{ch1994efimzelmanov:early-life}

Efim Zelmanov was born in Khabarovsk, in the Soviet Far East, in 1955.
He showed exceptional mathematical talent early and was drawn to
algebra.

\subsection{2.2 Education}\label{ch1994efimzelmanov:education}

Zelmanov studied at Novosibirsk State University, where he completed his
undergraduate and graduate work under Anatoly Shirshov. He received his
PhD in 1980 and his Doctor of Sciences degree in 1985.

\subsection{2.3 Academic Career}\label{ch1994efimzelmanov:academic-career}

\begin{itemize}
\tightlist
\item
 \textbf{1980--1986:} Researcher, Institute of Mathematics,
 Novosibirsk.
\item
 \textbf{1986--1987:} Professor, Novosibirsk State University.
\item
 \textbf{1987--1990:} Professor, University of Chicago.
\item
 \textbf{1990--1994:} Professor, University of Wisconsin, Madison.
\item
 \textbf{1994:} Fields Medal, International Congress of Mathematicians,
 Zürich.
\item
 \textbf{1994--1995:} Professor, Yale University.
\item
 \textbf{1995--2002:} Professor, University of Chicago.
\item
 \textbf{2002--present:} Professor, University of California, San
 Diego.
\item
 \textbf{2001:} Elected to the National Academy of Sciences.
\item
 \textbf{2012:} Elected to the American Academy of Arts and Sciences.
\end{itemize}

\begin{center}\rule{0.5\linewidth}{0.5pt}\end{center}

\section{3. The Restricted Burnside
Problem}\label{ch1994efimzelmanov:the-restricted-burnside-problem}

\subsection{3.1 The Burnside Problem}\label{ch1994efimzelmanov:the-burnside-problem}

The \textbf{Burnside problem} (posed by William Burnside in 1902) asks:
is every finitely generated group in which every element has finite
order necessarily finite? This was answered negatively by Golod and
Shafarevich in 1964, who constructed infinite groups where every element
has finite order.

\subsection{3.2 The Restricted Burnside
Problem}\label{ch1994efimzelmanov:the-restricted-burnside-problem-1}

The \textbf{restricted Burnside problem} is a variant: given integers
\(d\) and \(n\), is there a finite bound on the orders of all finite
\(d\)-generator groups of exponent \(n\)? In other words, is the
Burnside group \(B(d,n)\) finite? Here, the Burnside group is the
quotient of the free group on \(d\) generators by the normal subgroup
generated by all \(n\)-th powers.

\subsection{3.3 Historical Progress}\label{ch1994efimzelmanov:historical-progress}

The restricted Burnside problem had been solved for various small
exponents: - Exponent 2 (trivial --- all groups of exponent 2 are
abelian). - Exponent 3 (Burnside, 1902). - Exponent 4 (Sanov, 1940). -
Exponent 6 (Hall and Higman, 1956). For larger exponents, the problem
remained open for decades.

\subsection{3.4 Zelmanov's Solution}\label{ch1994efimzelmanov:zelmanovs-solution}

Zelmanov solved the restricted Burnside problem in full generality. The
proof is a masterpiece of algebra, combining group theory, Lie algebra
theory, and Jordan algebra theory.

\subsection{3.5 Significance}\label{ch1994efimzelmanov:significance}

The solution of the restricted Burnside problem is one of the great
achievements of 20th-century algebra. It completed a program of research
that had occupied some of the best mathematicians for nearly a century.

\begin{center}\rule{0.5\linewidth}{0.5pt}\end{center}

\section{4. Proof Structure: Groups to Lie
Algebras}\label{ch1994efimzelmanov:proof-structure-groups-to-lie-algebras}

\subsection{4.1 Reduction to Prime
Exponent}\label{ch1994efimzelmanov:reduction-to-prime-exponent}

The first step in the proof is to reduce the restricted Burnside problem
to the case of prime exponent. Using the Hall--Higman theorem and the
theory of finite groups, one can show that if the restricted Burnside
problem holds for all prime exponents, then it holds for all exponents.

\subsection{4.2 The Lie Algebra Associated to a
Group}\label{ch1994efimzelmanov:the-lie-algebra-associated-to-a-group}

For a group \(G\) of exponent \(p\), Zelmanov constructs a Lie algebra
\(\mathcal{L}(G)\) over \(\mathbb{F}_p\) using the descending central
series:

\[
\mathcal{L}(G) = \bigoplus_{i=1}^\infty \gamma_i(G)/\gamma_{i+1}(G)
\]

where \(\gamma_i(G)\) is the \(i\)-th term of the lower central series.
The Lie bracket is induced by the group commutator.

\subsection{4.3 The Engel Condition}\label{ch1994efimzelmanov:the-engel-condition}

Zelmanov shows that \(\mathcal{L}(G)\) satisfies the Engel condition:

\[
[x, y, y, \ldots, y] = 0
\]

for any \(x, y \in \mathcal{L}(G)\), where the number of \(y\)'s depends
only on \(p\).

\subsection{4.4 Local Nilpotency}\label{ch1994efimzelmanov:local-nilpotency}

The critical step is proving that Engel Lie algebras are locally
nilpotent: any finitely generated subalgebra is nilpotent. This step
uses the theory of Jordan algebras and sandwich algebras.

\subsection{4.5 Finiteness}\label{ch1994efimzelmanov:finiteness}

From the local nilpotency of \(\mathcal{L}(G)\), Zelmanov concludes that
\(G\) is locally finite, so the Burnside group \(B(d, n)\) is finite.

\begin{center}\rule{0.5\linewidth}{0.5pt}\end{center}

\section{5. Engel Lie Algebras}\label{ch1994efimzelmanov:engel-lie-algebras}

\subsection{5.1 Definition}\label{ch1994efimzelmanov:definition}

A Lie algebra \(L\) is called \textbf{Engel} if every element is
ad-nilpotent: for every \(x \in L\), the operator
\(\text{ad}_x(y) = [x, y]\) is nilpotent.

\subsection{5.2 The Zelmanov Theorem}\label{ch1994efimzelmanov:the-zelmanov-theorem}

Zelmanov proved that every Engel Lie algebra is \textbf{locally
nilpotent}: every finitely generated subalgebra is nilpotent. This is a
fundamental result in Lie algebra theory.

\subsection{5.3 The Role of Sandwich
Elements}\label{ch1994efimzelmanov:the-role-of-sandwich-elements}

Zelmanov discovered that Engel Lie algebras contain \textbf{sandwich
elements} --- non-zero elements \(a\) such that \([a, x, a] = 0\) for
all \(x\). The presence of these elements forces the algebra to be
nilpotent.

\begin{center}\rule{0.5\linewidth}{0.5pt}\end{center}

\section{6. Jordan Algebras}\label{ch1994efimzelmanov:jordan-algebras}

\subsection{6.1 Definition}\label{ch1994efimzelmanov:definition-1}

A \textbf{Jordan algebra} is a non-associative algebra satisfying: -
\(x \cdot y = y \cdot x\) (commutativity) -
\((x \cdot y) \cdot (x \cdot x) = x \cdot (y \cdot (x \cdot x))\)
(Jordan identity)

\subsection{6.2 Zelmanov's
Classification}\label{ch1994efimzelmanov:zelmanovs-classification}

Zelmanov classified all prime Jordan algebras, showing they fall into
four classes: - \textbf{Hermitian type:} The algebra of Hermitian
matrices over an associative algebra with involution. - \textbf{Clifford
type:} The algebra of a quadratic form. - \textbf{Albert type:} The
exceptional Albert algebra (the only exceptional simple Jordan algebra).
- \textbf{Quadratic form type:} The algebra of a bilinear form.

\subsection{6.3 The Zelmanov Prime
Theorem}\label{ch1994efimzelmanov:the-zelmanov-prime-theorem}

The \textbf{Zelmanov prime theorem} states that every prime Jordan
algebra is either of Hermitian, Clifford, Albert, or quadratic form
type. This is the Jordan analogue of the classification of prime
associative rings.

\subsection{6.4 Applications}\label{ch1994efimzelmanov:applications}

The classification has been used to: - Study the structure of
infinite-dimensional Jordan algebras. - Understand representations of
Jordan algebras. - Solve the Kurosh problem for Jordan algebras.

\begin{center}\rule{0.5\linewidth}{0.5pt}\end{center}

\section{7. The Kurosh Problem}\label{ch1994efimzelmanov:the-kurosh-problem}

\subsection{7.1 Statement}\label{ch1994efimzelmanov:statement}

The \textbf{Kurosh problem} (posed in 1941) asks: is every algebraic
algebra (an algebra where every element satisfies a polynomial identity)
locally finite-dimensional?

\subsection{7.2 Zelmanov's Solution}\label{ch1994efimzelmanov:zelmanovs-solution-1}

Zelmanov proved that every algebraic Jordan algebra is locally
finite-dimensional. This solved the Kurosh problem for Jordan algebras
and has implications for other classes of algebras.

\begin{center}\rule{0.5\linewidth}{0.5pt}\end{center}

\section{8. Lie Algebras in Positive
Characteristic}\label{ch1994efimzelmanov:lie-algebras-in-positive-characteristic}

\subsection{8.1 Background}\label{ch1994efimzelmanov:background}

Lie algebras over fields of characteristic \(p > 0\) behave very
differently from those in characteristic 0. New phenomena include: - The
presence of \(p\)-th power operations. - The existence of simple Lie
algebras that are not classical. - The relationship to group theory via
the restricted Burnside problem.

\subsection{8.2 Zelmanov's Contributions}\label{ch1994efimzelmanov:zelmanovs-contributions}

Zelmanov contributed to the theory of Lie algebras in positive
characteristic: - \textbf{The Engel theorem:} The generalization of
Engel's theorem to fields of characteristic \(p\). - \textbf{The
restricted Burnside problem:} Using Lie algebras in characteristic \(p\)
to study groups of exponent \(p\). - \textbf{Simple Lie algebras:} The
classification of simple Lie algebras over algebraically closed fields
of characteristic \(p\).

\begin{center}\rule{0.5\linewidth}{0.5pt}\end{center}

\section{9. The Structure of Nil
Algebras}\label{ch1994efimzelmanov:the-structure-of-nil-algebras}

\subsection{9.1 Definition}\label{ch1994efimzelmanov:definition-2}

A \textbf{nil algebra} is an algebra in which every element is
nilpotent. The Kurosh problem asks whether every nil algebra is locally
nilpotent.

\subsection{9.2 Zelmanov's Results}\label{ch1994efimzelmanov:zelmanovs-results}

Zelmanov proved: - Every nil Jordan algebra is locally nilpotent. -
Every nil Lie algebra satisfying the Engel condition is locally
nilpotent. - These results are key steps in the solution of the
restricted Burnside problem.

\begin{center}\rule{0.5\linewidth}{0.5pt}\end{center}

\section{10. The Theory of PI-Algebras}\label{ch1994efimzelmanov:the-theory-of-pi-algebras}

\subsection{10.1 Definition}\label{ch1994efimzelmanov:definition-3}

A \textbf{PI-algebra} (polynomial identity algebra) is an algebra that
satisfies a nontrivial polynomial identity. Examples include commutative
algebras, finite-dimensional algebras, and nilpotent algebras.

\subsection{10.2 Zelmanov's
Contributions}\label{ch1994efimzelmanov:zelmanovs-contributions-1}

Zelmanov contributed to the theory of PI-algebras: - The relationship
between PI-algebras and the Kurosh problem. - The structure of prime
PI-algebras. - The application of PI-theory to Jordan and Lie algebras.
- The use of PI-theory in the proof of the restricted Burnside problem.

\begin{center}\rule{0.5\linewidth}{0.5pt}\end{center}

\section{11. The Theory of Pro-p
Groups}\label{ch1994efimzelmanov:the-theory-of-pro-p-groups}

\subsection{11.1 Definition}\label{ch1994efimzelmanov:definition-4}

\textbf{Pro-p groups} are projective limits of finite p-groups. They
arise naturally in number theory as Galois groups of infinite
extensions.

\subsection{11.2 Zelmanov's
Contributions}\label{ch1994efimzelmanov:zelmanovs-contributions-2}

Zelmanov applied his work to pro-p groups: - \textbf{Groups of finite
exponent:} Pro-p groups of finite exponent are finite. - \textbf{The
Golod--Shafarevich theorem:} The structure of pro-p groups with few
relations. - \textbf{The relationship to the Burnside problem:} The
restricted Burnside problem is equivalent to the finiteness of certain
pro-p groups.

\begin{center}\rule{0.5\linewidth}{0.5pt}\end{center}

\section{12. Simple Groups and the Burnside
Problem}\label{ch1994efimzelmanov:simple-groups-and-the-burnside-problem}

\subsection{12.1 Background}\label{ch1994efimzelmanov:background-1}

The classification of finite simple groups is one of the great
achievements of 20th-century mathematics. The restricted Burnside
problem is closely related to the structure of finite simple groups.

\subsection{12.2 Zelmanov's
Contributions}\label{ch1994efimzelmanov:zelmanovs-contributions-3}

Zelmanov's solution implies: - Every finite simple group of exponent
\(n\) has order bounded by a function of \(n\). - The structure of
groups of exponent \(n\) is controlled by the exponent. - The
relationship between finite simple groups and Lie algebras.

\begin{center}\rule{0.5\linewidth}{0.5pt}\end{center}

\section{13. Graded Lie Algebras}\label{ch1994efimzelmanov:graded-lie-algebras}

\subsection{13.1 Definition}\label{ch1994efimzelmanov:definition-5}

A \textbf{graded Lie algebra} is a Lie algebra
\(L = \bigoplus_{i \in \mathbb{Z}} L_i\) such that
\([L_i, L_j] \subset L_{i+j}\). The lower central series of a group
gives a natural grading on the associated Lie algebra.

\subsection{13.2 Zelmanov's
Contributions}\label{ch1994efimzelmanov:zelmanovs-contributions-4}

Zelmanov studied graded Lie algebras arising from groups: - The grading
on \(\mathcal{L}(G)\) from the lower central series. - The Engel
condition in the graded setting. - The local nilpotency of graded Engel
Lie algebras.

\begin{center}\rule{0.5\linewidth}{0.5pt}\end{center}

\section{14. The Engel Condition: Deep
Structure}\label{ch1994efimzelmanov:the-engel-condition-deep-structure}

\subsection{14.1 Definition}\label{ch1994efimzelmanov:definition-6}

The \textbf{Engel condition} is the identity
\([x, y, y, \ldots, y] = 0\) where \(y\) appears \(n\) times. The
minimal such \(n\) is the Engel length.

\subsection{14.2 Zelmanov's Theorem}\label{ch1994efimzelmanov:zelmanovs-theorem}

Zelmanov proved that Engel Lie algebras are locally nilpotent. This
theorem is the key technical achievement and the heart of the Burnside
problem solution.

\subsection{14.3 The Proof Structure}\label{ch1994efimzelmanov:the-proof-structure}

The proof shows: 1. Engel algebras contain sandwich elements. 2. The set
of sandwich elements generates an ideal. 3. The quotient by this ideal
is again Engel. 4. By induction, the algebra is nilpotent.

\begin{center}\rule{0.5\linewidth}{0.5pt}\end{center}

\section{15. The Kostrikin Problem}\label{ch1994efimzelmanov:the-kostrikin-problem}

\subsection{15.1 Statement}\label{ch1994efimzelmanov:statement-1}

The \textbf{Kostrikin problem} asks: is every Engel Lie algebra of
bounded nilpotency index locally nilpotent? This is the Lie algebra
analogue of the restricted Burnside problem.

\subsection{15.2 Zelmanov's Solution}\label{ch1994efimzelmanov:zelmanovs-solution-2}

Zelmanov solved the Kostrikin problem as part of his work. His proof
shows that Engel Lie algebras with bounded Engel length are locally
nilpotent.

\begin{center}\rule{0.5\linewidth}{0.5pt}\end{center}

\section{16. The Hall--Higman Theorem}\label{ch1994efimzelmanov:the-hallhigman-theorem}

\subsection{16.1 Statement}\label{ch1994efimzelmanov:statement-2}

The \textbf{Hall--Higman theorem} is a fundamental result about
representations of finite solvable groups. It gives bounds on the
nilpotency class of certain Lie algebras.

\subsection{16.2 Role in Zelmanov's Work}\label{ch1994efimzelmanov:role-in-zelmanovs-work}

Zelmanov used the Hall--Higman theorem to: - Bound the Engel length of
\(\mathcal{L}(G)\) for groups of exponent \(p\). - Reduce the restricted
Burnside problem to prime exponents. - Control the structure of the Lie
algebra associated to a group.

\begin{center}\rule{0.5\linewidth}{0.5pt}\end{center}

\section{17. Free Groups and Free Burnside
Groups}\label{ch1994efimzelmanov:free-groups-and-free-burnside-groups}

\subsection{17.1 Free Burnside Groups}\label{ch1994efimzelmanov:free-burnside-groups}

A \textbf{free Burnside group} \(B(d, n)\) is the quotient of the free
group on \(d\) generators by the normal subgroup generated by all
\(n\)-th powers. Zelmanov's theorem shows that these groups are finite
for all \(d\) and \(n\).

\subsection{17.2 Structure of Free
Groups}\label{ch1994efimzelmanov:structure-of-free-groups}

Zelmanov studied: - The residual finiteness of free groups. - The
relationship between free groups and free Lie algebras. - The Magnus
theorem and its applications.

\begin{center}\rule{0.5\linewidth}{0.5pt}\end{center}

\section{18. The Theory of Sandwich
Algebras}\label{ch1994efimzelmanov:the-theory-of-sandwich-algebras}

\subsection{18.1 Definition}\label{ch1994efimzelmanov:definition-7}

A \textbf{sandwich algebra} is an algebra containing a non-zero element
\(a\) such that \(a x a = 0\) for all \(x\) in the algebra. Such
elements are called sandwich elements.

\subsection{18.2 Zelmanov's Discovery}\label{ch1994efimzelmanov:zelmanovs-discovery}

Zelmanov discovered that sandwich elements play a crucial role in Engel
Lie algebras. In any non-nilpotent Engel Lie algebra, there exists a
non-zero sandwich element. This leads to an inductive proof of
nilpotency.

\begin{center}\rule{0.5\linewidth}{0.5pt}\end{center}

\section{19. The Baer Radical}\label{ch1994efimzelmanov:the-baer-radical}

\subsection{19.1 Definition}\label{ch1994efimzelmanov:definition-8}

The \textbf{Baer radical} (or prime radical) of an associative ring is
the intersection of all prime ideals. For Lie algebras, it is defined
via the universal enveloping algebra.

\subsection{19.2 Zelmanov's
Contributions}\label{ch1994efimzelmanov:zelmanovs-contributions-5}

Zelmanov used the Baer radical to: - Study the structure of Engel Lie
algebras. - Relate the nilpotency of an algebra to its prime ideals. -
Connect the theory of Jordan algebras to Lie algebras.

\begin{center}\rule{0.5\linewidth}{0.5pt}\end{center}

\section{20. The Theory of Alternative
Algebras}\label{ch1994efimzelmanov:the-theory-of-alternative-algebras}

\subsection{20.1 Definition}\label{ch1994efimzelmanov:definition-9}

An \textbf{alternative algebra} is a non-associative algebra in which
the associator \((x, y, z) = (xy)z - x(yz)\) is alternating.

\subsection{20.2 Zelmanov's
Contributions}\label{ch1994efimzelmanov:zelmanovs-contributions-6}

Zelmanov contributed to the theory of alternative algebras: - The
relationship between alternative and Jordan algebras. - Prime
alternative algebras and their classification. - The structure of
alternative PI-algebras.

\begin{center}\rule{0.5\linewidth}{0.5pt}\end{center}

\section{21. The Theory of Varieties}\label{ch1994efimzelmanov:the-theory-of-varieties}

\subsection{21.1 Definition}\label{ch1994efimzelmanov:definition-10}

A \textbf{variety} of algebras is a class defined by a set of
identities. Varieties of groups, Lie algebras, and associative algebras
are fundamental objects in universal algebra.

\subsection{21.2 Zelmanov's
Contributions}\label{ch1994efimzelmanov:zelmanovs-contributions-7}

Zelmanov contributed to the theory of varieties: - The variety of groups
of exponent \(n\). - The variety of Engel Lie algebras. - The
relationship between varieties of groups and varieties of Lie algebras.

\begin{center}\rule{0.5\linewidth}{0.5pt}\end{center}

\section{22. The Golod--Shafarevich
Theorem}\label{ch1994efimzelmanov:the-golodshafarevich-theorem}

\subsection{22.1 Statement}\label{ch1994efimzelmanov:statement-3}

The \textbf{Golod--Shafarevich theorem} gives a sufficient condition for
a pro-p group to be infinite. It was used to construct the first
counterexamples to the Burnside problem.

\subsection{22.2 Zelmanov's
Contributions}\label{ch1994efimzelmanov:zelmanovs-contributions-8}

Zelmanov used the Golod--Shafarevich theorem to: - Understand which
Burnside groups are finite. - Study the structure of pro-p groups with
few relations. - Refine the Golod--Shafarevich inequality.

\begin{center}\rule{0.5\linewidth}{0.5pt}\end{center}

\section{23. The Novikov Problem}\label{ch1994efimzelmanov:the-novikov-problem}

\subsection{23.1 Statement}\label{ch1994efimzelmanov:statement-4}

The \textbf{Novikov problem} concerns the solvability of the word
problem for groups. The word problem asks whether there is an algorithm
to determine if two words represent the same element.

\subsection{23.2 Zelmanov's
Contributions}\label{ch1994efimzelmanov:zelmanovs-contributions-9}

Zelmanov's work on the restricted Burnside problem has implications for
the Novikov problem: - The finiteness of Burnside groups implies that
certain group-theoretic problems are decidable. - The relationship
between the word problem and the structure of finite groups.

\begin{center}\rule{0.5\linewidth}{0.5pt}\end{center}

\section{24. The Chernikov Problem}\label{ch1994efimzelmanov:the-chernikov-problem}

\subsection{24.1 Statement}\label{ch1994efimzelmanov:statement-5}

The \textbf{Chernikov problem} asks whether a group in which all
subgroups are finite-by-abelian is itself finite-by-abelian.

\subsection{24.2 Zelmanov's
Contributions}\label{ch1994efimzelmanov:zelmanovs-contributions-10}

Zelmanov applied his methods to the Chernikov problem: - Using
techniques from the restricted Burnside problem. - Studying groups with
finiteness conditions. - The structure of infinite groups with finite
subgroups.

\begin{center}\rule{0.5\linewidth}{0.5pt}\end{center}

\section{25. Growth in Groups}\label{ch1994efimzelmanov:growth-in-groups}

\subsection{25.1 Definition}\label{ch1994efimzelmanov:definition-11}

The \textbf{growth of a group} measures the number of group elements
expressible as words of length at most \(n\) in the generators. Groups
can have polynomial, exponential, or intermediate growth.

\subsection{25.2 Zelmanov's
Contributions}\label{ch1994efimzelmanov:zelmanovs-contributions-11}

Zelmanov's work on the restricted Burnside problem has implications for
growth: - Burnside groups have bounded growth. - The relationship
between exponent and growth rate. - The existence of groups of
intermediate growth.

\begin{center}\rule{0.5\linewidth}{0.5pt}\end{center}

\section{26. The Hilbert Nullstellensatz and
PI-Algebras}\label{ch1994efimzelmanov:the-hilbert-nullstellensatz-and-pi-algebras}

\subsection{26.1 The Nullstellensatz}\label{ch1994efimzelmanov:the-nullstellensatz}

The \textbf{Hilbert Nullstellensatz} is a fundamental result in
algebraic geometry relating polynomial equations to algebraic varieties.

\subsection{26.2 Zelmanov's
Contributions}\label{ch1994efimzelmanov:zelmanovs-contributions-12}

Zelmanov's work connects to algebraic geometry through: - The
Nullstellensatz for PI-algebras. - The relationship between the Burnside
problem and the theory of algebraic varieties. - The use of
combinatorial methods in algebraic geometry.

\begin{center}\rule{0.5\linewidth}{0.5pt}\end{center}

\section{27. The Magnus Theorem}\label{ch1994efimzelmanov:the-magnus-theorem}

\subsection{27.1 Statement}\label{ch1994efimzelmanov:statement-6}

The \textbf{Magnus theorem} states that free groups are residually
finite: for any non-trivial element, there is a finite quotient in which
the element remains non-trivial.

\subsection{27.2 Zelmanov's Uses}\label{ch1994efimzelmanov:zelmanovs-uses}

Zelmanov used the Magnus theorem in his work on the restricted Burnside
problem, as residual finiteness is essential for understanding the
structure of Burnside groups.

\begin{center}\rule{0.5\linewidth}{0.5pt}\end{center}

\section{28. Engel Identities in Associative
Rings}\label{ch1994efimzelmanov:engel-identities-in-associative-rings}

\subsection{28.1 Engel Rings}\label{ch1994efimzelmanov:engel-rings}

An \textbf{Engel ring} is an associative ring satisfying
\([x, y, y, \ldots, y] = 0\) where \([x, y] = xy - yx\).

\subsection{28.2 Zelmanov's
Contributions}\label{ch1994efimzelmanov:zelmanovs-contributions-13}

Zelmanov studied Engel rings: - Engel rings are locally nilpotent under
certain conditions. - Engel rings give rise to Engel Lie algebras. - The
connection between Engel rings and the restricted Burnside problem.

\begin{center}\rule{0.5\linewidth}{0.5pt}\end{center}

\section{29. Awards and Recognition}\label{ch1994efimzelmanov:awards-and-recognition}

\textbf{1994:} - Fields Medal (International Mathematical Union)

\textbf{1996:} - Elected to the Russian Academy of Sciences

\textbf{2001:} - Elected to the National Academy of Sciences (USA)

\textbf{2012:} - Elected to the American Academy of Arts and Sciences

\textbf{1996:} - Humboldt Research Award

\begin{center}\rule{0.5\linewidth}{0.5pt}\end{center}

\section{30. Selected Major
Publications}\label{ch1994efimzelmanov:selected-major-publications}

\begin{enumerate}
\def\labelenumi{\arabic{enumi}.}
\item
 \textbf{``The restricted Burnside problem.''} \emph{Proceedings of the
 International Congress of Mathematicians}, Zürich, 1994. --- \emph{The
 solution of the restricted Burnside problem.}
\item
 \textbf{``Engel Lie algebras and the restricted Burnside problem.''}
 \emph{Mathematics of the USSR-Sbornik} 71, 467--478, 1992. ---
 \emph{The key theorem on Engel Lie algebras.}
\item
 \textbf{``Prime Jordan algebras.''} \emph{Algebra i Logika} 18,
 147--179, 1979. --- \emph{The classification of prime Jordan
 algebras.}
\item
 \textbf{``The Kurosh problem for Jordan algebras.''} \emph{Siberian
 Mathematical Journal} 24, 1983. --- \emph{The solution of the Kurosh
 problem for Jordan algebras.}
\item
 \textbf{``Absolute zero divisors in Jordan algebras.''}
 \emph{Izvestiya Akademii Nauk SSSR} 45, 1981. --- \emph{The theory of
 sandwich algebras.}
\item
 \textbf{``Nilpotency of Engel Lie algebras.''} \emph{Communications in
 Algebra} 19, 1991. --- \emph{The local nilpotency of Engel Lie
 algebras.}
\item
 \textbf{``The Golod--Shafarevich theorem and the restricted Burnside
 problem.''} \emph{Proceedings of the Steklov Institute} 234, 2001. ---
 \emph{The relationship between the Golod--Shafarevich theorem and the
 Burnside problem.}
\end{enumerate}

\begin{center}\rule{0.5\linewidth}{0.5pt}\end{center}

\section{31. References and Further
Reading}\label{ch1994efimzelmanov:references-and-further-reading}

\subsection{Official Sources}\label{ch1994efimzelmanov:official-sources}

\begin{itemize}
\tightlist
\item
 IMU Fields Medal 1994 Announcement:
 https://www.mathunion.org/imu-awards/fields-medal/fields-medal-1994
\item
 UCSD: ``Zelmanov awarded Fields Medal.'' 1994.
\end{itemize}

\subsection{Key Technical References}\label{ch1994efimzelmanov:key-technical-references}

\begin{itemize}
\tightlist
\item
 Zelmanov, E. ``The restricted Burnside problem.'' ICM 1994.
\item
 Zelmanov, E. ``Engel Lie algebras.'' \emph{Mathematics of the
 USSR-Sbornik} 71, 1992.
\item
 Zelmanov, E. ``Prime Jordan algebras.'' \emph{Algebra i Logika} 18,
 1979.
\end{itemize}

\subsection{Biography and Overview}\label{ch1994efimzelmanov:biography-and-overview}

\begin{itemize}
\tightlist
\item
 Kostrikin, A. ``The Work of Efim Zelmanov.'' ICM 1994 Laudatio.
\item
 Vaughan-Lee, M. ``Efim Zelmanov and the restricted Burnside problem.''
 \emph{Notices of the AMS}, 1995.
\item
 Jacobson, N. \emph{Jordan Algebras}. Springer, 1968.
\end{itemize}

\subsection{Textbooks}\label{ch1994efimzelmanov:textbooks}

\begin{itemize}
\tightlist
\item
 Vaughan-Lee, M. \emph{The Restricted Burnside Problem}. Oxford
 University Press, 1990.
\item
 Jacobson, N. \emph{Structure and Representations of Jordan Algebras}.
 AMS, 1968.
\item
 Kostrikin, A. \emph{Around Burnside}. Springer, 1990.
\end{itemize}

\begin{center}\rule{0.5\linewidth}{0.5pt}\end{center}

\emph{(End of report)}
